
For downstream usability, event-level $F_1$ did not by itself establish suitability for blink-rate or PERCLOS-style measures, because these applications require predicted counts to vary with the true counts across recordings~\citep{dai2023detection,zulkarnanie2022enhancements,ren2019comparison,alyan2023blink}. Proposed-Med provided the closest count fidelity, with a mean predicted-to-true ratio of 1.157, Pearson $r=0.941$, and Lin CCC$=0.938$. BLINKER-concat instead produced approximately twice the true count, with a ratio of 1.939, despite retaining a relatively high correlation of $r=0.868$; this pattern indicates systematic inflation rather than accurate absolute measurement. MNE-annot showed the converse limitation: its ratio of 1.036 was near unity, but its lower correlation of $r=0.629$ and CCC$=0.617$ indicate that over-counting and under-counting offset in the pooled mean rather than that session-level counts were faithfully tracked. Proposed-Med therefore better preserves between-session variation needed for downstream ocular indices. However, its positive bias of +80.2 events and ratio above unity show that it still over-counted, so absolute blink rates and derived thresholds should be interpreted with care.
